% !PRINTED 527
\weberpage{527}
Nun ist nach einem bekannten Satz aus der Theorie der
$\Gamma$-Funktionen, wenn $k$ eine beliebige Größe mit positiv imaginärem
 Bestandteil bedeutet:
\[ \frac{1}{(-ik)^s} = \frac{(2\pi)^s}{\Gamma(s)}\int_0^\infty e^{2\pi ik\xi}\,\xi^{s-1}\,d\xi, \]
worin die Potenz $(-ik)^s$ dadurch eindeutig erklärt ist, daß, wenn
$-ik = \varrho e^{i\Theta}$ gesetzt, $\varrho$ positiv und der Winkel $\Theta$ zwischen $-\frac{\pi}{2}$
und $+\frac{\pi}{2}$ angenommen wird, $(-ik)^s = \varrho^s e^{i\Theta s}$ zu setzen ist.

Ersetzt man hierin $ik$ durch die konjugiert imaginäre
Größe $-ik'$, indem man zugleich einen neuen Integrationsbuchstaben
 $\eta$ wählt, und multipliziert beide Gleichungen, so folgt:
\weq{7}{\frac{1}{(kk')^s} = \frac{(2\pi)^{2s}}{\Gamma(s)\Gamma(s)}\int_0^\infty\int_0^\infty e^{2\pi i(k\xi-k'\eta)}(\xi\eta)^{s-1}\,d\xi\,d\eta.}
Hierin setzen wir für $k, k'$ die beiden konjugierten Faktoren
$x+\omega_1y$, $x-\omega_2y$ von $\psi$ und führen den Integralausdruck (7)
in (6) ein.

Dadurch erhalten wir:
\weq{8}{M_s = \frac{(2\pi)^{2s}}{A_s\Gamma(s)\Gamma(s)}\Sigma^{x,y}2\int_0^\infty\int_0^\infty e^{2\pi i[x(\xi-\eta)+y(\omega_1\xi+\omega_2\eta)]}(\xi\eta)^{s-1}\,d\xi\,d\eta.}
% SOURCE-ORIGINAL-SIC: Equation (8) prints A_s in the denominator although equation (6) and the derivation indicate A^s; the source subscript is preserved.
Wir fassen nun das Produkt der beiden Integrale auf der
rechten Seite dieses Ausdruckes, das wir zur Abkürzung durch
\weq{9}{W = 2\int_0^\infty\int_0^\infty e^{2\pi i[x(\xi-\eta)+y(\omega_1\xi+\omega_2\eta)]}(\xi\eta)^{s-1}\,d\xi\,d\eta}
bezeichnen, als Doppelintegral auf, das sich, wenn $\xi, \eta$ als rechtwinkelige
 Koordinaten in der Ebene gedeutet werden, über den positiven
Quadranten des Koordinatensystems erstreckt.

% Fig. 2 source description: axes eta and xi; the bisector separates regions I and II. At fixed u the hatched integration strips extend from v=u toward v=infinity, not from the origin.
\weberfigureimage{weber_b3_i0558_fig2_source_crop.png}{0.30}

Um das Doppelintegral umzuformen, teilen wir den Integrationsbereich
 durch eine den Winkel halbierende Gerade in die beiden
Teile I, II (s. die Fig. 2), und substituieren im ersten Teil
%CONT>>
